% =========================================================
%  SECCIÓN VI — CONCLUSIONES
% =========================================================
\section{Conclusiones}

Este trabajo present\'o el an\'alisis comparativo de cuatro
estrategias de vectorizaci\'on SIMD aplicadas al kernel num\'erico
del operador Prewitt para detecci\'on de bordes en im\'agenes.
A partir de los resultados obtenidos se establecen las siguientes
conclusiones:

\begin{enumerate}
  \item \textbf{El kernel Prewitt es memory-bound.}
    Su intensidad aritm\'etica te\'orica de $0.333$~FLOP/Byte
    lo coloca a la izquierda del punto de cresta del modelo
    Roofline ($0.387$~FLOP/Byte), lo que significa que el
    cuello de botella de rendimiento es el ancho de banda de
    memoria y no la capacidad de c\'omputo del procesador.
    Esta conclusi\'on es independiente del nivel de vectorizaci\'on
    aplicado y fue confirmada por las cuatro gr\'aficas Roofline
    generadas con Intel Advisor.

  \item \textbf{La vectorizaci\'on produce ganancias modestas.}
    El speedup m\'aximo observado fue de $1.24\times$ para la
    vectorizaci\'on autom\'atica y de $1.21\times$ para la
    expl\'icita.
    En un kernel \textit{compute-bound}, estos mismos anchos
    vectoriales (16$\times$ y 32$\times$ efectivo) producir\'ian
    speedups proporcionales al ancho SIMD.
    La modestia de los resultados es, por tanto, esperada y
    correcta dado el car\'acter \textit{memory-bound} del kernel.

  \item \textbf{La vectorizaci\'on autom\'atica es la m\'as eficiente.}
    El compilador \texttt{icpx} con \texttt{-O3 -march=native}
    elige un ancho vectorial de~16~floats y genera instrucciones
    AVX2 \'optimas sin intervenci\'on del programador, superando
    incluso a la versi\'on expl\'icita en t\'erminos de relaci\'on
    rendimiento-esfuerzo.
    Esto subraya el alto nivel de madurez de los optimizadores
    modernos de compiladores Intel.

  \item \textbf{La vectorizaci\'on guiada requiere ajuste del simdlen.}
    El pragma \texttt{simdlen(8)} impuso un ancho menor al
    elegido autom\'aticamente por el compilador, resultando en
    rendimiento pr\'acticamente igual al escalar.
    Una elecci\'on de \texttt{simdlen(16)} o dejando que el
    compilador decida habr\'ia producido mejores resultados.

  \item \textbf{El 88\% del programa es c\'odigo escalar.}
    La fracci\'on vectorizable del programa (el kernel) representa
    solo el 11--12\% del tiempo total de ejecuci\'on.
    El tiempo dominante corresponde a operaciones inherentemente
    escalares: lectura/escritura de archivos PGM, c\'alculo
    de gradientes 2D con m\'ascaras 3$\times$3, y replicaci\'on
    de datos en memoria.
    Esta distribuci\'on limita el speedup global del programa de
    acuerdo con la Ley de Amdahl.

  \item \textbf{Intel Advisor es una herramienta diagn\'ostica
    valiosa.}
    El Survey, los trip counts y el modelo Roofline permitieron
    confirmar objetivamente el ancho vectorial real de cada
    versi\'on, identificar el cuello de botella de memoria y
    cuantificar la fracci\'on vectorizable del programa.
    Su integraci\'on en el flujo de desarrollo es recomendable
    para cualquier proyecto de optimizaci\'on de kernels HPC.
\end{enumerate}

\section*{Agradecimientos}

El autor agradece al Dr.~Mario Rossainz L\'opez, profesor del
curso de Programaci\'on en Plataformas Multicore de la Benem\'erita
Universidad Aut\'onoma de Puebla (BUAP), por el dise\~no de la
actividad y la gu\'ia metodol\'ogica que permitieron llevar a cabo
este trabajo de manera rig\'urosa.
